\subsection{INDUCED FILTER}

PROPOSITION 8. Let \(\mathfrak{F}\) be a filter on a set X and A a subset of X. Then the trace \(\mathfrak{F}_{\mathrm{A}}\) of \(\mathfrak{F}\) on A is a filter if and only if each set of \(\mathfrak{F}\) meets A.

Since \((\mathbf{M} \cap \mathbf{N}) \cap \mathbf{A} = (\mathbf{M} \cap \mathbf{A}) \cap (\mathbf{N} \cap \mathbf{A})\) we see that \(\mathfrak{F}_{\mathbf{A}}\) satisfies \((\mathbf{F}_{\mathbf{II}})\) ; again, if \(\mathbf{M} \cap \mathbf{A} \subset \mathbf{P} \subset \mathbf{A}\) then \(\mathbf{P} = (\mathbf{M} \cup \mathbf{P}) \cap \mathbf{A}\) , whence \(\mathfrak{F}_{\mathbf{A}}\) satisfies \((\mathbf{F}_{\mathbf{I}})\) . Hence \(\mathfrak{F}_{\mathbf{A}}\) is a filter if and only if it satisfies \((\mathbf{F}_{\mathbf{III}})\) , i.e.~if and only if each set of \(\mathfrak{F}\) meets \(\mathbf{A}\) .

In particular, if \(\mathbf{A} \in \mathfrak{F}\) then \(\mathfrak{F}_{\mathbf{A}}\) is a filter on \(\mathbf{A}\) , by \((\mathbf{F}_{\mathbf{II}})\) and \((\mathbf{F}_{\mathbf{III}})\) .

DEFINITION 5. Let A be a subset of a set X and F a filter on X. If the trace of F on A is a filter on A, this filter is said to be induced by F on A.

If a filter \(\mathfrak{F}\) on \(X\) induces a filter on \(A \subset X\) , then the trace on \(A\) of a base of \(\mathfrak{F}\) is a base of \(\mathfrak{F}_A\) , by reason of Proposition 3 of no. 3.

Example. Let X be a topological space, A a subset of X, x a point of X. In order that the trace on A of the neighbourhood filter B of x should be a filter on A, it is necessary and sufficient that every neighbourhood of x meets A, i.e.~that x lies in the closure of A (§ 1, no. 6, Definition 10).

This example of an induced filter is of interest for two reasons: first because it plays an important role in the theory of limits (§ 7, no. 5) and secondly because every filter can be defined in this way. Indeed, let F be a filter on a set X and let X' be the set obtained by adjoining a new element ω to X, X being identified with the complement of \{ω\} in X' (Set Theory, R, § 4, no. 5); let F' be the filter on X' consisting of the sets M ∪ \{ω\} where M runs through F. For each point x ≠ ω of X', let B(x) be the set of all subsets of X' which contain x, and let B(ω) be F'; then the B(x) for x ∈ X' obviously satisfy axioms (V₁), (V₂), (V₃) and (V₄) and therefore define a topology on X' for which they are the neighbourhood filters of points. Finally ω lies in the closure of X in this topology, and F is induced by F' = B(ω) on X. The topology thus defined on X' (resp. the set X' with this topology) is called the topology (resp. the topological space) associated with F.

PROPOSITION 9. An ultrafilter \(\mathfrak{U}\) on a set \(X\) induces a filter on a subset \(A\) of \(X\) if and only if \(A \in \mathfrak{U}\) ; and if this condition is satisfied then \(\mathfrak{U}_A\) is an ultrafilter on \(A\) .

This is an immediate consequence of Propositions 5 and 6 of no. 4.

